\section{Results}
\label{sec:results}

\subsection{Overall Angle-Slot Regression}
\label{sec:overall_results}

Teacher-forced regression covers 521 angle slots from 310 held-out geometry-grounded records.
YogaLLaVA V3 obtains an MAE of $13.98^\circ$ and an RMSE of $21.82^\circ$, with a Pearson correlation of 0.913.
The test $R^2$ is 0.832 and differs from the validation $R^2$ of 0.874 used for checkpoint selection (Sec.~\ref{sec:implementation_details}).
In addition, 57.01\% and 70.63\% of predictions fall within $10^\circ$ and $15^\circ$, respectively, of the reconstruction-derived references.
Table~\ref{tab:overall_results} in Appendix~\ref{app:additional_results} lists the full metric set.
These results indicate that the model can identify and express the requested measurement from the supplied geometry under a known answer structure.

Two reference points bound the teacher-forced result.
The first is a constant predictor that returns the test-set mean at every slot.
Such a predictor attains an $R^2$ of 0 by construction, and its RMSE equals the standard deviation of the targets.
The reported $R^2$ of 0.832 and RMSE of $21.82^\circ$ place that standard deviation at $53.24^\circ$.
YogaLLaVA V3 therefore reduces RMSE by 59\% relative to the constant predictor.
The second is an oracle that reads the requested descriptor out of the geometry input and copies its value.
Because 518 of the 521 targets are supplied among the geometry tokens (Sec.~\ref{sec:evaluation_protocol}), the oracle attains an MAE of $0^\circ$ and an $R^2$ of 1 on those slots.
The distance from the oracle measures how often YogaLLaVA V3 selects the wrong descriptor or fails to reproduce a value already supplied to it.


\subsection{Analysis by Question Type}
\label{sec:path_results}

Table~\ref{tab:path_results} reports performance across the five geometry-QA families.
Misconception questions achieve the lowest MAE ($11.28^\circ$) and highest within-$10^\circ$ accuracy (68.75\%), although this subset contains only 32 slots.
Regional-summary and comparison questions also perform favorably despite requiring multiple values or cross-body-part relations.
Instructor-observation questions are the most challenging, with an MAE of $21.08^\circ$ and within-$10^\circ$ accuracy of 25.81\%.
Their $R^2$ of 0.054 over 31 slots is the weakest result in the table and indicates that predictions in this family barely track their targets.
The MAE and threshold accuracies alone understate the weakness.
Their less explicit teacher-style formulations may make descriptor selection more ambiguous, although the present experiment does not isolate this factor.

\begin{table}[t]
    \centering
    \caption{Teacher-forced regression by geometry-QA type.}
    \label{tab:path_results}
    \resizebox{\columnwidth}{!}{
    \begin{tabular}{lrrrrr}
        \toprule
        \textbf{Question Type} & \textbf{Slots} &
        \textbf{RMSE $\downarrow$} & \textbf{MAE $\downarrow$} &
        \textbf{$R^2$ $\uparrow$} & \textbf{$\leq10^\circ$ $\uparrow$} \\
        \midrule
        Comparison & 106 & $21.53^\circ$ & $12.95^\circ$ & 0.791 & 62.26\% \\
        Direct angle & 155 & $22.65^\circ$ & $15.22^\circ$ & 0.821 & 51.61\% \\
        Instructor observation & 31 & $28.69^\circ$ & $21.08^\circ$ & 0.054 & 25.81\% \\
        Misconception & 32 & $17.10^\circ$ & $11.28^\circ$ & 0.877 & 68.75\% \\
        Region summary & 197 & $20.72^\circ$ & $12.88^\circ$ & 0.847 & 61.42\% \\
        \bottomrule
    \end{tabular}}
\end{table}

\subsection{Analysis by Geometric Descriptor}
\label{sec:metric_results}

Among descriptors with substantial test support, trunk twist yields the lowest MAE ($6.84^\circ$), followed by lateral trunk bend, trunk tilt, and right shoulder elevation (Table~\ref{tab:metric_results} in Appendix~\ref{app:additional_results}).
Left shoulder elevation is the most difficult descriptor, followed by elbow flexion and right hip opening.
These differences may arise from descriptor-specific distributions, reconstruction uncertainty, occlusion, or left--right ambiguity.
Identifying their individual effects requires further controlled analysis.

The symmetry categories contain only three or four slots and should not be used to draw broad conclusions.
Hip symmetry is also distinct because no dedicated hip-symmetry token is included in the 14-descriptor input vocabulary.
Its three test targets therefore cannot be retrieved from a corresponding independent input token.
Despite the missing token, hip symmetry reaches an MAE of $7.82^\circ$, lower than several descriptors with dedicated input tokens and far larger test support.
Under a retrieval-based reading of the regression, the direction of this result is counterintuitive.
We leave the discrepancy as an open question given the three-slot sample.


\subsection{Autoregressive Generation}
\label{sec:generation_comparison}

Table~\ref{tab:generation_results} evaluates free-running generation.
YogaLLaVA V3 produces the correct number of numeric values for 267 of 310 records (86.1\%), compared with 66.1\% for LLaVA-1.5-7B and 68.1\% for Qwen2.5-VL-7B.
The Wilson 95\% confidence interval for YogaLLaVA V3 is $[81.8,89.5]\%$ and overlaps neither baseline interval ($[60.7,71.2]\%$ for LLaVA-1.5-7B and $[62.7,73.0]\%$ for Qwen2.5-VL-7B).
YogaLLaVA V3 is therefore more reliable at determining when and how many continuous outputs are required.

This comparison evaluates complete models.
It does not isolate a single architectural factor.
YogaLLaVA V3 receives the geometry input and task-specific instruction tuning, whereas the baselines are image-only and zero-shot.
The observed advantage reflects the combined effect of geometry conditioning, domain adaptation, and continuous angle regression.
The comparison should not be read as a modality-matched ablation.

For numeric comparison, we use the 160-record common-success subset on which all three models produce the correct slot count.
Conditioning on joint success removes every record where a baseline produced an invalid slot count.
Because such records are plausibly the harder ones, the numeric gaps below hold only over the records that all three models handled.
YogaLLaVA V3 achieves a slot-level MAE of $17.65^\circ$, compared with $42.80^\circ$ for LLaVA and $46.91^\circ$ for Qwen.
Its record-level macro MAE is $16.98^\circ$ with a 95\% bootstrap confidence interval of $[14.15,20.17]^\circ$.
The corresponding intervals are $[37.33,46.41]^\circ$ for LLaVA and $[41.59,53.87]^\circ$ for Qwen.
Paired bootstrap differences relative to YogaLLaVA are $24.85^\circ$ ($[19.45,30.35]^\circ$) for LLaVA and $30.64^\circ$ ($[23.70,37.83]^\circ$) for Qwen.

\begin{table}[t]
    \centering
    \caption{Autoregressive generation on 310 test records. Slot accuracy is
    measured on the full test set and reported with Wilson 95\% confidence
    intervals. Numeric errors use the 160-record common-success subset.}
    \label{tab:generation_results}
    \resizebox{\columnwidth}{!}{
    \begin{tabular}{lrrrr}
        \toprule
        \textbf{Model} & \textbf{Slot Match $\uparrow$} &
        \textbf{95\% CI} &
        \textbf{MAE $\downarrow$} & \textbf{RMSE $\downarrow$} \\
        \midrule
        YogaLLaVA V3 & 267/310 (86.1\%) & $[81.8,89.5]\%$ & $17.65^\circ$ & $27.23^\circ$ \\
        LLaVA-1.5-7B & 205/310 (66.1\%) & $[60.7,71.2]\%$ & $42.80^\circ$ & $52.59^\circ$ \\
        Qwen2.5-VL-7B & 211/310 (68.1\%) & $[62.7,73.0]\%$ & $46.91^\circ$ & $61.12^\circ$ \\
        \bottomrule
    \end{tabular}}
\end{table}

\subsection{Discussion and Limitations}
\label{sec:result_interpretation}

The teacher-forced and autoregressive evaluations characterize complementary capabilities.
The former measures continuous prediction at known answer slots, whereas the latter additionally requires the model to decide whether a numeric slot is needed, produce and order the correct number of slots, preserve consistency between numbers and language, and terminate appropriately.
Reporting them separately avoids conflating regression accuracy with generation quality.

The current evidence supports geometry-conditioned measurement selection and expression.
The references inherit uncertainty from monocular reconstruction, including viewpoint, occlusion, foreshortening, clothing, inversion, depth, and left--right ambiguities.
SAM 3D Body reports a mean per-joint position error of $54.8$~mm on 3DPW~\cite{yang2026sam3dbody}.
That figure is positional and does not convert to an angular tolerance.
No joint-angle error is published for the reconstruction, and we therefore cannot separate the share of the reported $13.98^\circ$ MAE that originates in the references themselves.
The references should not be interpreted as clinical, diagnostic, rehabilitation, motion-capture, or manually verified biomechanical ground truth.

The misconception and instructor-observation families place a further limitation on the present scope.
Both train the model to issue corrective, teacher-style judgments, and their supervision comes from the same measurements we decline to certify as clinically meaningful.
We have not checked the resulting corrections against expert instruction.
The generated judgments are therefore unvalidated as instructional advice.

The domain-knowledge set carries a separate caveat.
It supplies 4,090 of the 10,090 records and takes part in training, yet no experiment here measures whether the model answers questions about benefits, contraindications, or alignment cues correctly.
We therefore claim it as an annotation resource and make no accuracy claim for pose-level knowledge.

Finally, the experiments do not isolate modality contributions through ablation and do not test generalization to unseen yoga-pose categories.
These gaps, together with domain-knowledge evaluation, remain directions for future work.
